\chapter{INTRODUCTION}

In today’s healthcare environment, efficient management of patient records, appointments, prescriptions, laboratory tests, and billing processes is essential for providing quality medical services. Traditional hospital management systems that rely heavily on manual documentation and paper-based records often result in data redundancy, delayed communication, increased human error, and difficulties in maintaining patient history. To overcome these limitations, healthcare institutions require an integrated and automated digital solution.

The Medora Hospital Management System is a web-based application developed to streamline and automate hospital operations through a centralized and secure platform. The system integrates multiple hospital departments—including Administration, Reception, Doctors, Pharmacy, and Laboratory—into a unified digital ecosystem. It ensures smooth data flow between departments, reduces manual workload, and enhances operational transparency.

The system implements role-based access control, allowing different users such as Admin, Receptionist, Doctor, Patient, Pharmacist, and Lab Staff to perform tasks according to their responsibilities. The Reception module handles patient registration, appointment scheduling, and billing processes. Doctors manage consultations, generate digital prescriptions, and request laboratory tests. Lab staff update and upload test reports electronically, while pharmacists manage medicine inventory and dispensing.

In addition to consultation-based treatment, the system provides flexibility by allowing patients to register and purchase medicines directly from the pharmacy without mandatory doctor consultation. Patients can also upload external prescriptions, which are verified and processed by pharmacists. This feature improves accessibility while maintaining proper medication tracking and record management.

Developed using the Laravel framework and MySQL database, the Medora Hospital Management System ensures data security, scalability, and performance efficiency. By digitizing patient records and automating hospital workflows, the system enhances administrative control, improves communication between departments, and supports modern healthcare service delivery.
